\section{Introduction}
Spatial prediction models are often used to estimate responses at locations where observations are unavailable. Assessing their prediction error requires a validation procedure that represents the intended prediction task. Cross-validation (CV) estimates error by holding out groups of observations, fitting the prediction model to the remaining data, and comparing predictions with the held-out responses. However, the spatial relationships created by these splits can differ from those encountered when the final prediction model is applied. With clustered sampling, validation observations may have close training neighbours while much of the prediction area lies farther from the data. Conversely, holding out large spatial blocks can create greater separations than the final prediction task requires. These differences can affect how well CV estimates target prediction error.

Nearest-neighbour distance matching uses the prediction locations to guide the choice of validation training sets. \citet{mila2022nndm} proposed NNDM leave-one-out CV, which selectively excludes nearby training observations to match validation and target nearest-distance distributions over distances within an autocorrelation range. \citet{linnenbrink2024knndm} extended this approach to $k$-fold CV through kNNDM. It compares candidate fold configurations using the Wasserstein distance $W$, an established measure of the difference between probability distributions \citep{panaretos2019wasserstein}. Here, $W$ measures the total area between the cumulative distributions of validation and target nearest-training distances. Choosing a smaller $W$ brings these two distributions closer together.

The usefulness of nearby training observations also depends on the prediction model and its predictors. Spatial prediction error profiles estimate how a particular model's observed error changes with distance from training data \citep{brenning2023profiles}. They can reveal whether increasing separation is associated with increasing error for that model. Geographic proximity alone also cannot establish whether predictor values at a target location are represented in training. Area-of-applicability methods assess this aspect of prediction by comparing target predictor values with those in the training data \citep{meyer2021aoa}.

Other spatial validation methods adjust CV scores to account for correlation or differences between the sampled and target locations. Covariance-based corrections account for correlated responses \citep{rabinowicz2022correlated,yuval2025classification}. Sampling-intensity weights, spatial modelling of validation residuals, and importance weights combined with exclusion buffers address differences between the sampled sites and the prediction target \citep{debruin2022clustered,wang2023geospatial}. Dissimilarity-adaptive CV combines random and spatial CV scores according to predictor similarity \citep{wang2025dissimilarity}, while target-weighted CV weights validation losses to represent target characteristics \citep{brenning2026targetweighted}.

For distance matching, the central question is how a difference in nearest-training distance translates into a difference in prediction error. The Wasserstein distance describes the geographic difference without using the strength of correlation at each separation. For a prediction model that uses spatial dependence, increasing separation matters when it removes useful information from neighbouring responses. If responses are already weakly correlated at both distances, a larger separation may have little effect. Where correlation decreases rapidly, a smaller change in distance may remove more information. This suggests evaluating distance matching through its relationship with spatial correlation and the prediction model's mean squared error (MSE).

Our first contribution is an explicit connection between nearest-neighbour correlation and MSE. For Gaussian prediction using one neighbour, with known mean and covariance, the reduction in MSE is proportional to the squared correlation with that observation. This motivates $J_2$, which subtracts the average squared correlation at target locations from the corresponding average during validation. Using standard Gaussian conditioning \citep{rasmussen2006gpml}, we show that response variance multiplied by $J_2$ is exactly target MSE minus expected CV MSE.

Our second contribution is to quantify what a nearest-neighbour comparison leaves out. When Gaussian prediction uses all training observations, the remaining neighbours can reduce error beyond what the nearest observation alone achieves. We derive how this additional reduction differs between validation and target locations. This difference can have either sign and forms the part of expected CV bias that $J_2$ does not capture. A leave-one-out counterexample gives identical nearest-distance distributions but nonzero expected CV bias. The result explains why even perfect nearest-distance matching may be insufficient when prediction uses several observations.

Our third contribution is to test whether selecting CV schemes using correlation gives more accurate estimates of target prediction error. We compare selection by the smallest $|J_2|$ with selection by the smallest $W$ among random folds, spatial blocks, buffered spatial blocks and kNNDM. Studies with known and estimated covariance assess expected bias and the accuracy of the observed CV error estimate. Further simulations vary the prediction model and the information supplied by predictors while holding response covariance fixed. An application predicts continuous soil-cadmium concentrations in the Swiss Jura and compares the selected CV estimates with measured prediction error at reserved sites.

Section~\ref{sec:methods} defines the prediction target and diagnostics, derives their relationships with MSE, and describes the CV selection rules. Section~\ref{sec:simulation} presents the simulation designs and results. Section~\ref{sec:application} presents the Swiss Jura application and its findings. Section~\ref{sec:discussion} discusses the implications for choosing and interpreting spatial CV schemes.
